\chapter{Conclusion and Future Work}

\setlength{\parindent}{1em}

\section{Conclusion}

This thesis set out to test a single hypothesis: that the dominant obstacle to
consistent personal-finance tracking is data-entry friction, and that
conversational bilingual speech, processed by a self-hosted, open-source AI
pipeline, can remove most of that friction without sacrificing privacy or
incurring per-request cost. MiPay realizes and substantiates that hypothesis as a
complete, reproducible system.

The work delivered, against the objectives set out in Chapter~1, the following.
It designed and implemented an end-to-end pipeline that converts a bilingual
voice or typed utterance into one or more structured, validated transaction
records (O1). It selected and justified an open ASR model (Whisper via
\texttt{faster-whisper}) and an open instruction-tuned language model
(Qwen2.5-7B-Instruct via Ollama) that run on commodity CPU-only hardware (O2). It
guaranteed machine-parseable extraction through JSON-schema constrained decoding,
eliminating format-level failure by construction (O3). It implemented a
deterministic, unit-tested post-processing layer---numeral normalization,
dialect-aware relative-date resolution, currency defaulting, and confidence-based
review routing---that supplies the exactness the language model cannot (O4). And
it built a fully localized Arabic/English mobile client with three input modes and
a human-in-the-loop confirmation workflow, alongside a labeled bilingual
evaluation corpus and a reproducible harness reporting Word Error Rate and
per-field extraction accuracy with ablation studies (O5, O6).

The central intellectual contribution is not any single component but their
\emph{composition}: a robust multilingual speech model, a few-shot
schema-constrained language model that needs no task-specific training data, and
a deterministic exact-computation layer, joined by a stateless, human-confirmed
save path. This architecture is what lets a self-hosted system that is smaller
and cheaper than a commercial cloud service nonetheless be a usable product,
because the human confirmation step converts the model's occasional errors into
one-tap corrections rather than corrupted data. The self-hosted choice, finally,
is what gives the project its academic substance---explicit model selection,
measurable metrics, and principled ablations---rather than the thin-wrapper
character of a cloud-API integration.

\section{Limitations}

The project is honest about what it does not yet do. The evaluation corpus,
though spanning Arabic, English, and code-switched input, is authored by the team
and may not match the distribution of real user speech; the audio subset used for
end-to-end evaluation is smaller than the text corpus. Voice processing is
synchronous and, on CPU-only hardware, takes several seconds per clip---adequate
for interactive single-user use but not for high concurrency. The monthly summary
sums amounts without currency conversion, reporting in the user's default
currency. The extraction model, being prompted rather than fine-tuned, occasionally
mislabels currency or category on rare phrasings, and dialect coverage in the
date resolver, while broad, is not exhaustive. The system requires connectivity to
the backend for voice processing; there is no offline mode. None of these
limitations undermines the core contribution, but each defines a direction for
future work.

\section{Future Work}

The future-work program follows directly from the project's design philosophy:
keep the system self-hosted and privacy-preserving, and improve each stage with
newer or task-specialized models as resources allow. The current
implementation---Whisper \texttt{small} plus Qwen2.5-7B---was selected as the best
quality/cost balance on the CPU-only reference machine; the upgrades below are the
principled next steps, each stated with its rationale and its expected trade-off.

\subsection{Algorithmic Enhancements}

\paragraph{Upgraded speech model.} The most direct quality gain is to replace
Whisper \texttt{small} with \texttt{faster-whisper large-v3} or its distilled
\texttt{large-v3-turbo} variant. These substantially reduce Arabic and
code-switched WER while, in the \texttt{turbo} case, retaining most of the speed
advantage of CTranslate2 \texttt{int8} inference~\cite{fasterwhisper,ctranslate2}.
The trade-off is memory and latency: \texttt{large-v3} is practical with a GPU
(NFR-01's $\leq$\,4\,s budget) or on a higher-RAM CPU host. The pipeline already
exposes the model size as a configuration variable, so this is a deployment
change, not a code change.

\paragraph{Upgraded extraction model.} The extractor can be moved to a newer or
larger open-weight model---a Qwen3-class model or a 14B-parameter instruct model---to
improve currency and category disambiguation and multi-transaction splitting on
rare phrasings, while preserving the constrained-decoding guarantee that already
makes output schema-valid. Because the extraction interface depends only on the
Ollama \texttt{format} contract, swapping the model is again a configuration
change. The trade-off is inference cost; 4-bit quantization~\cite{frantar2023gptq}
keeps even a 14B model within the reference memory envelope at some latency cost.

\paragraph{Task-specific fine-tuning.} The natural academic extension is to stop
prompting and start learning. Every confirmed (user-corrected) transaction is a
gold label pairing a transcript with its correct fields; accumulating these yields
exactly the labeled, in-domain, bilingual corpus that the supervised approaches of
Chapter~2 require but that does not exist today. A parameter-efficient LoRA
fine-tune~\cite{hu2022lora} of a small model---or a fine-tuned multilingual NER
head~\cite{antoun2020arabert,lample2016ner}---trained on this corpus could match or
exceed the prompted 7B model at a fraction of the inference cost, closing the loop
between product usage and model improvement.

\subsection{System Improvements}

\paragraph{Asynchronous processing.} For higher concurrency, the synchronous
voice endpoint can be moved behind a job queue (FastAPI background tasks or
Celery) with a polling or push-based result endpoint, decoupling request latency
from inference time. The pipeline is already structured to make this a
localized change.

\paragraph{On-device inference.} A fully offline mode---\texttt{whisper.cpp} plus a
1--3B quantized language model running on the handset---would remove the backend
dependency entirely for privacy-maximalist users, at the cost of reduced accuracy
and increased battery use.

\paragraph{Streaming UX.} Partial transcripts streamed during recording, and
progressive extraction results, would make the interaction feel instantaneous even
when total processing takes several seconds.

\subsection{Feature Expansion}

Beyond the AI pipeline, the product roadmap includes budgets and savings goals,
automatic recurring-transaction detection, receipt OCR and bank-SMS parsing as
additional structured-extraction input modes (natural extensions of the same
schema-constrained approach), multi-currency conversion in the dashboard, and data
export. Each reuses the existing confirm-and-save core.

\subsection{Evaluation Framework}

The evaluation itself is a direction for growth. The most immediate item is
re-running the full evaluation against the current twenty-nine-category,
EGP-default schema (Section~4.5.1): the corpus and results reported in
Chapter~5 were measured against the earlier seventeen-category configuration,
and confirming that the two root causes identified there (Arabic
compound-number parsing, prompt-coverage gaps) hold at the same or greater
severity with twelve additional categories is the natural next data point.
Beyond that: expanding the corpus with real, user-corrected utterances;
reporting \emph{per-dialect} WER and F\textsubscript{1} (Gulf, Egyptian,
Levantine) rather than a single Arabic figure; and adding a longitudinal study
of how often users must edit extractions in real use---the truest measure of
the friction reduction the project set out to achieve.

\section{Concluding Remarks}

MiPay demonstrates that the components needed for a private, low-cost, bilingual
voice interface to a real application are now available as open, self-hostable
building blocks, and that assembling them well---with constrained decoding for
safety, a deterministic layer for exactness, and a human checkpoint for
trust---produces a system that is both academically rigorous and practically
usable. The models will improve, and the documented upgrade path makes adopting
those improvements a matter of configuration rather than redesign; but the
architecture---fuzzy understanding by the model, exact computation by code,
final judgment by the user---is the durable contribution of this work.
